\begin{abstract}
Large language models (LLMs) deliver strong reasoning and generation performance, but their cost, latency, energy demand, and hardware requirements remain high. Small language models (SLMs) are better suited to edge and privacy-sensitive deployment, yet their capabilities vary across domains. This paper presents a systematic literature review of 54 studies on collaborative SLM--LLM systems spanning hybrid, multi-agent, retrieval-centered, and ensemble architectures, with monolithic LLMs treated as baselines rather than as a primary architecture class. Guided by four research questions on performance, efficiency, orchestration, and domain suitability, we apply a Kitchenham-guided review protocol, using PRISMA 2020 terminology to report the selection flow, together with a three-pass qualitative coding scheme that synthesizes models, architectures, datasets, metrics, latency, cost, energy, and reported limitations. Corpus analysis shows that 72.2\% of studies date to 2025 and that hybrid orchestration is the most frequently tagged design under the multi-label architecture coding; a contribution-type mapping confirms that solution proposals dominate. Efficiency reporting is uneven across dimensions: latency is reported in roughly one-third of the corpus and cost in fewer than one-quarter, whereas direct energy measurement appears in only a single study, so cost and energy conclusions cannot be treated as equally evidenced. Outcome coding indicates that strong results track how well an explicit control layer (router, retriever, planner) is calibrated to the task rather than model scale; because the eligibility criteria required every included study to employ some control mechanism, the corpus cannot contrast the presence of such a layer against its absence. Overall, SLMs are most competitive when they are paired with well-calibrated routing, retrieval, planning, or fusion mechanisms in narrow, structured, or domain-specific settings.
\end{abstract}
